% ╔══════════════════════════════════════════════════════════════╗
% ║          package_and_settings.tex — Template_Slides          ║
% ║  Beamer 簡報模板的套件載入與全域設定                          ║
% ║  （本檔不需要每次改，切換文件只改 main.tex 的兩行 \input）    ║
% ╚══════════════════════════════════════════════════════════════╝


% ── 1. 色彩主題選擇 ──────────────────────────────────────────────
%
%  在這裡更改一個字，整份簡報的配色就會跟著換。
%  可選值：green（預設，深青綠色調）、blue（深藍色調）、purple（紫色調）
%
\def\slidetheme{orange}   % ← 改這裡切換色調（用 \def，不用 \newcommand）


% ── 依選擇套用對應色系（\ifx 比較展開後字串）────────────────────
%  注意：比較用的 \themeoptiongreen / \themeoptionblue 名稱不含 @，
%        避免在 .tex 檔中 @ 被當成非字母而解析錯誤。
\def\themeoptiongreen{green}
\def\themeoptionblue{blue}
\def\themeoptionorange{orange}

\ifx\slidetheme\themeoptiongreen
  % ── 綠色系（預設）────────────────────────────────────────────
  \definecolor{darktitle}   {RGB}{30, 110,  70}  % 標題列底色：深綠
  \definecolor{accentcolor} {RGB}{30, 110,  70}  % 強調色：深綠
  \definecolor{bodytextcolor}{RGB}{24,  88,  56}  % 正文字色：深綠（同解答色）
  \definecolor{answerblue}  {RGB}{24,  88,  56}  % 解答字色：深綠
  \definecolor{alertcolor}  {RGB}{200,  0,   0}  % 警示色：紅
  \definecolor{examplefill} {RGB}{232, 246, 236} % problembox 底色：淡綠
  \definecolor{exampleline} {RGB}{60,  140,  80} % problembox 框線：中綠
  \definecolor{solutionfill}{RGB}{240, 248, 243} % 投影片底色
\else\ifx\slidetheme\themeoptionblue
  % ── 藍色系 ───────────────────────────────────────────────────
  \definecolor{darktitle}   {RGB}{25,  55, 120}  % 標題列底色：深藍
  \definecolor{accentcolor} {RGB}{25,  55, 120}  % 強調色：亮藍
  \definecolor{bodytextcolor}{RGB}{20,  44,  96}  % 正文字色：深藍（同解答色）
  \definecolor{answerblue}  {RGB}{20,  44,  96}  % 解答字色：深藍
  \definecolor{alertcolor}  {RGB}{200,  0,   0}  % 警示色：紅
  \definecolor{examplefill} {RGB}{230, 238, 252} % problembox 底色：淡藍
  \definecolor{exampleline} {RGB}{50,  110, 200} % problembox 框線：亮藍
  \definecolor{solutionfill}{RGB}{249, 251, 254} % 投影片底色
\else\ifx\slidetheme\themeoptionorange
  % ── 橘色系（UT Austin Burnt Orange）──────────────────────────
  \definecolor{darktitle}   {RGB}{191,  87,   0}  % 標題列底色：Burnt Orange
  \definecolor{accentcolor} {RGB}{191,  87,   0}  % 強調色：Burnt Orange
  \definecolor{bodytextcolor}{RGB}{153,  70,   0}  % 正文字色：深橘（同解答色）
  \definecolor{answerblue}  {RGB}{153,  70,   0}  % 解答字色：深橘
  \definecolor{alertcolor}  {RGB}{200,   0,   0}  % 警示色：紅
  \definecolor{examplefill} {RGB}{253, 240, 228} % problembox 底色：淡橘
  \definecolor{exampleline} {RGB}{191,  87,   0} % problembox 框線：Burnt Orange
  \definecolor{solutionfill}{RGB}{251, 244, 236} % 投影片底色
\else
  % ── 紫色系 ───────────────────────────────────────────────────
  \definecolor{darktitle}   {RGB}{62,  30, 105}  % 標題列底色：深紫
  \definecolor{accentcolor} {RGB}{62,  30, 105}  % 強調色：深紫
  \definecolor{bodytextcolor}{RGB}{50,  24,  84}  % 正文字色：深紫（同解答色）
  \definecolor{answerblue}  {RGB}{50,  24,  84}  % 解答字色：深紫
  \definecolor{alertcolor}  {RGB}{200,   0,   0} % 警示色：紅
  \definecolor{examplefill} {RGB}{240, 234, 250} % problembox 底色：淡紫
  \definecolor{exampleline} {RGB}{62,  30, 105}  % problembox 框線：深紫
  \definecolor{solutionfill}{RGB}{244, 240, 251} % 投影片底色
\fi\fi\fi

% 輔助顏色（跨主題共用）
\colorlet{lightbg}{solutionfill}  % 投影片底色
\definecolor{mygreen    }{RGB}{0,  153,  76}  % \textcolor{mygreen}{} 用
\definecolor{myblue     }{RGB}{0,  102, 204}  % \textcolor{myblue}{}  用


% ── 2. 字型設定（XeLaTeX 必用）─────────────────────────────────
\usepackage{amsmath}                % 數學公式核心（需在 unicode-math 之前）
\usepackage{mathtools}
\usepackage{amssymb}
\PassOptionsToPackage{no-math}{fontspec}  % 避免 fontspec 覆蓋數學字型
\usepackage{unicode-math}           % 數學字型（需在 fontspec 之前）
\setmathfont[
  Path     = ./fonts/,
  BoldFont = {KpMath-Bold.otf}
]{KpMath-Regular.otf}

\usepackage{xeCJK}                  % 中文字型支援

% 西文主字型：LinLibertine（正文），Kp Roman（粗體/斜體）
\setmainfont[
  Path            = ./fonts/,
  Numbers         = OldStyle,
  ItalicFont      = {KpRoman-Italic.otf},
  BoldFont        = {KpRoman-Bold.otf},
  BoldItalicFont  = {KpRoman-BoldItalic.otf}
]{LinLibertine_R.otf}

% 中文主字型：圓體
\setCJKmainfont[Path=./fonts/, BoldFont={cwTeXQMing-Medium.ttf}, BoldFeatures={FakeBold=1.5}]{cwTeXQMing-Medium.ttf}
\setCJKfamilyfont{round}[Path=./fonts/, BoldFont={cwTeXQYuan-Medium.ttf}, BoldFeatures={FakeBold=1.5}]{cwTeXQYuan-Medium.ttf}

% 中文額外字族與對應指令
\setCJKfamilyfont{kai} [Path=./fonts/]{cwTeXQKai-Medium.ttf}
\newcommand{\kai} [1]{{\CJKfamily{kai}#1}}   % 楷體

\setCJKfamilyfont{ming}[Path=./fonts/]{cwTeXQMing-Medium.ttf}
\newcommand{\ming}[1]{{\CJKfamily{ming}#1}}  % 明體（\yuan 同效果）
\newcommand{\yuan}[1]{{\CJKfamily{ming}#1}}  % 舊名保留

\setCJKfamilyfont{hei} [Path=./fonts/]{cwTeXQHei-Bold.ttf}
\newcommand{\hei} [1]{{\CJKfamily{hei}#1}}   % 黑體
\newcommand{\round}[1]{{\CJKfamily{round}#1}} % 圓體


% ── 3. 數學與科學套件 ────────────────────────────────────────────
\usepackage{bm}                     % \bm{} 粗斜體向量符號
\DeclareMathOperator{\Var}{Var}     % 自訂算子：Var
\DeclareMathOperator{\Cov}{Cov}     % 自訂算子：Cov
\DeclareMathOperator{\E}{\mathbb{E}} % 自訂算子：期望值


% ── 4. 排版與工具套件 ────────────────────────────────────────────
\usepackage{setspace}               % 行距控制
\onehalfspacing                     % 1.5 倍行距

\usepackage{booktabs}               % 表格專用橫線（\toprule \midrule \bottomrule）
\usepackage{graphicx}               % 插圖（\includegraphics）

\usepackage{tikz}                   % 向量圖形繪製
\usetikzlibrary{shapes, arrows, positioning}  % TikZ 流程圖元件

\usepackage{tcolorbox}              % 自訂色彩框（用於 problembox 等環境）
\tcbuselibrary{skins, breakable}    % breakable：允許框跨多頁（Beamer 需用 allowframebreaks）

\usepackage{hyperref}               % 超連結
\hypersetup{
  colorlinks = true,
  linkcolor  = black,
  filecolor  = magenta,
  urlcolor   = accentcolor,         % URL 顏色跟隨主題
}
\urlstyle{same}


% ── 5. Beamer 主題設定 ───────────────────────────────────────────
\usetheme{Boadilla}                 % 主題：Boadilla（簡約風格）
\setbeamerfont{title}{series=\mdseries, family=\CJKfamily{round}}  % 封面標題不用粗體
\setbeamerfont{author}{family=\CJKfamily{round}}
\setbeamerfont{date}{family=\CJKfamily{round}}
\setbeamerfont{frametitle}{series=\bfseries, family=\CJKfamily{round}}  % frame 標題加粗

% 套用色系到 Beamer 各元件
%  ► frametitle 底色 = darktitle（深色），文字白色
%  ► structure / item / section in toc 使用 accentcolor，跟隨主題變色
%    green → 深綠   blue → 深藍   purple → 淡薰衣草紫
%  ► 正文 (normal text) 維持黑色，確保可讀性
\setbeamercolor{frametitle}            {bg=darktitle,    fg=white}
\setbeamercolor{normal text}           {bg=lightbg,      fg=bodytextcolor}
\setbeamercolor{item}                  {fg=accentcolor}
\setbeamercolor{title}                 {fg=bodytextcolor}
\setbeamercolor{structure}             {fg=accentcolor}
\setbeamercolor{section in toc}        {fg=accentcolor}
\setbeamercolor{alerted text}          {fg=alertcolor}

% 項目符號樣式
\setbeamertemplate{itemize items}[circle]

% 關閉右下角 Beamer 預設導航圖示（通常用不到）
\beamertemplatenavigationsymbolsempty


% ── 6. 頁尾（Footer）設定 ────────────────────────────────────────
%
%  頁尾三欄：[目前章節] [簡報標題] [頁碼 / 總頁數]
%  頁尾底色配合主題自動更換。
%
\setbeamercolor{section in head/foot}     {bg=darktitle!10, fg=bodytextcolor}
\setbeamercolor{title in head/foot}       {bg=darktitle!30, fg=bodytextcolor}
\setbeamerfont{section in head/foot}      {family=\CJKfamily{round}}
\setbeamerfont{title in head/foot}        {family=\CJKfamily{round}}
\setbeamercolor{author in head/foot}      {bg=darktitle,    fg=white}
\setbeamercolor{page number in head/foot} {bg=darktitle,    fg=lightbg}
\setbeamercolor{subsection in head/foot}  {bg=gray!20,      fg=white}
\setbeamercolor{date in head/foot}        {bg=gray!10,      fg=black}

\setbeamertemplate{footline}{%
  \leavevmode%
  \hbox{%
    \begin{beamercolorbox}[wd=.33\paperwidth, ht=2.5ex, dp=1ex, center]
      {section in head/foot}%
      \insertsection
    \end{beamercolorbox}%
    \begin{beamercolorbox}[wd=.34\paperwidth, ht=2.5ex, dp=1ex, center]
      {title in head/foot}%
      \insertshorttitle
    \end{beamercolorbox}%
    \begin{beamercolorbox}[wd=.33\paperwidth, ht=2.5ex, dp=1ex, right]
      {page number in head/foot}%
      \insertframenumber{} / \inserttotalframenumber\hspace*{2ex}
    \end{beamercolorbox}%
  }%
  \vskip0pt%
}


% ── 7. 自訂環境 ──────────────────────────────────────────────────
%
%  Beamer 環境說明：
%  ┌─────────────────────────────────────────────────────┐
%  │  一個 frame 對應一張投影片。如果內容超過一頁，         │
%  │  請在 \begin{frame} 加上 [allowframebreaks] 選項：    │
%  │    \begin{frame}[allowframebreaks]{標題}              │
%  │  這樣 Beamer 會自動把超出的內容延伸到下一張投影片。   │
%  └─────────────────────────────────────────────────────┘

% ── problembox：題目框（有色框線，底色淡）──────────────────────
%
%  用法：\begin{problembox}{1.1}{題目名稱} ... \end{problembox}
%
%  ► 若題目很長、需要跨多頁，請在外層 frame 加 [allowframebreaks]
%
\newtcolorbox{problembox}[2]{
  enhanced,
  breakable,                          % 允許跨頁（需 frame [allowframebreaks]）
  colback  = white,                   % 框內白底
  colframe = exampleline,             % 框線色（隨主題變色）
  fonttitle= \bfseries,
  title    = {\round{題目 #1 \hspace{0.5em} #2}},
  before upper = {},
}

% ── answer：解答（藍字，不加框）─────────────────────────────────
%
%  用法：\begin{answer} ... \end{answer}
%
\newenvironment{answer}{\par\leftskip=1.5em\color{answerblue}}{\par}

% ── highlightbox：一般強調框（可自訂標題）──────────────────────
%
%  用法：\begin{highlightbox}{框標題} ... \end{highlightbox}
%
\newtcolorbox{highlightbox}[1]{
  enhanced,
  breakable,
  colback  = white,
  colframe = accentcolor,
  fonttitle = \bfseries,
  title    = {\round{#1}},
}

% ── Beamer 內建 block 的顏色也跟著主題走 ───────────────────────
\setbeamercolor{block title}         {bg=darktitle,   fg=white}
\setbeamercolor{block body}          {bg=examplefill, fg=black}
\setbeamercolor{block title example} {bg=accentcolor, fg=white}
\setbeamercolor{block body example}  {bg=examplefill, fg=black}
\setbeamercolor{block title alerted} {bg=alertcolor,  fg=white}
\setbeamercolor{block body alerted}  {bg=examplefill, fg=black}
